\section{Reduction to the section construction}\label{sec:completion}

The sole input from the remaining sections is
\Cref{prop:section-construction}. We first give the reduction in every
sign case. This identifies exactly what the geometric construction must
produce, including its assertion at a prescribed point of the original
open base.

\subsection{Conventions and birational changes}
All varieties are over $\C$, except when a generic fiber is explicitly
considered over a characteristic-zero function field.  We use additive
notation for line bundles and rational Cartier divisors.  A statement about
sections of a rational divisor is understood after taking a positive
Cartier multiple.  For numerical divisor classes, $P\ge_{\mathrm{pe}}Q$
means that $P-Q$ is pseudoeffective.

We make projective resolutions and resolve rational maps throughout.
Resolution may be taken to preserve a specified smooth open set and to be
equivariant for a finite group; see
\cite[Theorems 3.35 and 3.36]{KollarResolution}.
Here is the boundary convention that accompanies a birational
modification.  If $\mu:Z'\to Z$ is a birational morphism between smooth
projective varieties and $\Delta$ is reduced SNC, put
\[
 \Delta'=\mu^{-1}_*\Delta+\operatorname{Exc}(\mu)_{\mathrm{red}},
\]
after further resolution when necessary.  The log canonicity of the
smooth pair $(Z,\Delta)$ gives
\begin{equation}\label{eq:prelim-birational-log}
 K_{Z'}+\Delta'=\mu^*(K_Z+\Delta)+A_\mu,
 \qquad A_\mu\geq0\quad\text{$\mu$-exceptional}.
\end{equation}
Consequently pullback identifies the spaces of logarithmic pluricanonical
sections in every degree.  Indeed, pullback supplies one inclusion; in the
other direction a rational section regular away from a set of codimension
at least two on the smooth variety $Z$ extends across that set.  Equivalently,
$\mu_*\OO_{Z'}(mA_\mu)=\OO_Z$.  This convention also applies to a generic
fiber after a characteristic-zero extension of the ground field.

We use the standard numerical properties of big and nef divisors
\cite[Chapters 1--2]{Lazarsfeld}.  In particular, the big cone is the
interior of the pseudoeffective cone, and adding a pseudoeffective class
to a big class preserves bigness.  If $P$ is nef on an $n$-dimensional
projective variety, its numerical dimension is
\begin{equation}\label{eq:prelim-numerical-dimension}
 \nu(P)=\max\{k\in\{0,\ldots,n\}:P^k A^{n-k}>0\},
 \qquad A\ \text{ample},
\end{equation}
independently of $A$, and $P$ is big exactly when $\nu(P)=n$.
Bigness is preserved in both directions by a dominant generically finite
map.  For descent, factor the map through its finite normalization and
apply the field norm to an effective multiple of the pulled-back divisor
minus a sufficiently small pullback of an ample divisor.  Birational
invariance then gives the assertion on a resolution.  On a point the empty
intersection product is $1$, and numerical and Iitaka dimensions of the
zero line are both zero.

\subsection{The two negative-infinity cases}

\phantomsection\label{proof:fiber-vanishing}
Suppose first that $\kappa(F,L_F)=-\infty$. A nonzero section of $mL_X$
cannot vanish identically on every very general fiber of a dominant map:
its nonvanishing locus is a nonempty open subset of $X$, whose image
contains a dense open subset of $Y$. Restriction and adjunction would
therefore give a nonzero section of $mL_F$ on a very general fiber. This is
impossible. Hence $\kappa(X,L_X)=-\infty$.

\phantomsection\label{proof:base-vanishing}
Suppose next that $\kappa(Y,L_Y)=-\infty$. A hypothetical nonzero section
of $mL_X$ contradicts the first assertion of
\Cref{prop:section-construction}.
Thus the total Kodaira dimension is again $-\infty$. This is an actual
vanishing conclusion, not a consequence of an automatic lower inequality.

\subsection{A base of Kodaira dimension zero}

Assume $\kappa(Y,L_Y)=0$. If $Y$ is a point, the conclusion is immediate.
Otherwise choose $y\in V$ very general so that all fiber pluriform spaces
have their very general dimensions and $y$ lies outside the zero divisor
of every nonzero base pluriform. This last condition is possible:
\begin{equation}\label{eq:base-one-section}
 h^0(Y,aL_Y)\leq 1\qquad(a>0).
\end{equation}
Indeed, two linearly independent sections in a common degree have a
nonconstant ratio and hence give Iitaka dimension at least one. There are
only countably many degrees to consider.

For every $m>0$, restriction is injective:
\begin{equation}\label{eq:fiber-injective}
 H^0(X,mL_X)\longrightarrow H^0(F,mL_F).
\end{equation}
To prove this, suppose that a nonzero section on the left vanishes on $F$.
By \Cref{prop:section-construction}, this gives a nonzero section of a
positive multiple of $L_Y$ vanishing at $y$. That proposition applies at
every point of the original $V$, regardless of the auxiliary open sets
used for this particular section. This contradicts the choice of
$y$. Thus \eqref{eq:fiber-injective} holds in every degree.

The usual growth characterization of Iitaka dimension, applied to
\eqref{eq:fiber-injective}, gives
\begin{equation}\label{eq:zero-base-bound}
 \kappa(X,L_X)\leq\kappa(F,L_F).
\end{equation}
One may use sufficiently divisible degrees on both sides; the elementary
growth statement does not require finite generation of either section
ring, cf.\ \cite[Section~2.1]{Lazarsfeld}.

\subsection{Positive and intermediate base dimensions}

Let $k=\kappa(Y,L_Y)>0$. We use the reduction through the base Iitaka
fibration of Popa--Schnell \cite[\S2, Step~2]{PopaSchnellSmooth}, retaining
the reduced boundary and its strata. Resolve a logarithmic Iitaka fibration
$q\colon Y'\to T$, with $Y',T$ smooth projective and connected fibers, so
that $\dim T=k$. Put $D'$ equal to the strict transform of $D$ plus the
reduced exceptional divisor, with further resolution to SNC. The section
spaces of $K_{Y'}+D'$ agree with those of $L_Y$.

The very general fiber $G$ of $q$, with $D_G=D'|_G$, has
\begin{equation}\label{eq:iitaka-fiber-zero}
 \kappa(G,K_G+D_G)=0.
\end{equation}
Here the Iitaka map is taken with the relative algebraic closure of its
image field. The standard section-ratio argument proving
\eqref{eq:iitaka-fiber-zero} is the same one used in
\Cref{lem:relative-iitaka-construction}: nonzero forms restrict to the general fiber,
and an additional independent ratio there would spread after clearing
vertical poles with a sufficiently positive twist from the image. The
resolved original system dominates that twist, contradicting maximality
of the Iitaka image. No abundance statement is being used.

Take the main-component birational base change of $X$ to $Y'$ and resolve
it compatibly with the boundary. Denote the resulting morphism by
$f'\colon X'\to Y'$ and its boundary by $E'$, the strict transform of $E$
plus the reduced exceptional divisor. We can preserve the smooth pair
over $Y'\setminus D'$; the support inclusion holds because over $V$ the
pullback of a center of codimension at least two still has codimension at
least two, whereas the original inverse image of $D$ already lies in the
boundary. Thus all necessary new divisors are exceptional on the source
or lie over the old boundary. The morphism has connected fibers: it has
geometrically connected generic fiber, and its Stein factor is finite
birational over the normal base. These modifications preserve
$\kappa(X,L_X)$ and $\kappa(Y,L_Y)$.

For a very general point $t\in T$, put $G=q^{-1}(t)$ and
$H=(q\circ f')^{-1}(t)$. By generic smoothness applied to the finitely
many boundary strata, these are smooth projective varieties with their
induced reduced SNC boundaries.  Locally, surjectivity on the deepest
stratum lets the boundary parameters and the parameters pulled back from
$T$ be part of one coordinate system; restriction therefore retains
distinct reduced normal-crossing boundary branches.  The induced map
\[
 f_t\colon (H,E'|_H)\longrightarrow(G,D'|_G)
\]
has connected fibers and satisfies the original stratum-smoothness
hypothesis over $G\setminus D'|_G$. The condition over that open is
preserved by base change. Adjunction gives the restrictions of the two
log canonical divisors. The very general fiber invariant of $f_t$ is
$\kappa(F,L_F)$: choose $t$ and then the point of $G$ outside the
countably many loci governing all section dimensions. Applying
\eqref{eq:zero-base-bound}, or the already proved negative-infinity fiber
case, yields
\[
 \kappa(H,K_H+E'|_H)\leq\kappa(F,L_F).
\]
For completeness, take any nonempty total pluriform system, with
rational image map $\phi$. The image of the combined map
$(q\circ f',\phi)$ has generic fiber over $T$ equal to the image of the
restricted subsystem. Its fiber dimension is at most
$\kappa(H,K_H+E'|_H)$. Projection to the image of $\phi$ therefore gives
$\dim\operatorname{im}\phi\leq\dim T+\kappa(H,K_H+E'|_H)$.
Taking the supremum over degrees gives
\[
 \kappa(X',K_{X'}+E')
 \leq \dim T+\kappa(H,K_H+E'|_H)
 \leq k+\kappa(F,L_F).
\]
If the middle fiber Kodaira dimension is $-\infty$, restriction excludes
all total sections instead. This proves \Cref{thm:upper-bound}, including
the case $k=\dim Y$ and hence a point Iitaka fiber.
